\subsection*{最初に必要な用語}
\addcontentsline{toc}{subsection}{最初に必要な用語}

\par\smallskip\noindent\refstepcounter{table}\label{tab:ch04-01}表 \thetable: 最初に必要な用語の整理\par\smallskip
表 \thetable は、最初に必要な用語の整理を比較・確認しやすい形で整理したものである。\par\smallskip
\begin{longtable}{L{0.23\linewidth}L{0.69\linewidth}}
\toprule
用語 & 初学者向けの説明 \\
\midrule
ソフトウェア構築 & ソフトウェアを詳細に作成・保守する活動である。コーディング、検証、単体テスト、統合テスト、デバッグなどを含む。 \\
\term{コーディング}{Coding} & プログラミング言語などを使い、設計や要求を実行可能な形に書き表すこと。 \\
検証 & 作ったものが、仕様、設計、標準、期待された振る舞いを満たしているか確認すること。 \\
\term{単体テスト}{Unit Testing} & クラス、関数、モジュールなど、個別の部品を単独で確認するテスト。 \\
統合テスト & 複数の部品を組み合わせたとき、相互作用が正しく働くかを確認するテスト。 \\
\term{デバッグ}{Debugging} & 失敗の原因となる欠陥を見つけ、原因を理解し、修正する作業。 \\
依存関係 & 自分のソフトウェアが利用する外部ライブラリ、フレームワーク、部品、サービスなどとの関係。 \\
\term{継続的統合}{Continuous Integration} & 変更を頻繁に統合し、自動ビルドと自動テストで早く問題を発見する実践。 \\
再利用 & 既存の部品、コード、ライブラリ、サービス、設計パターンなどを別の問題解決に使うこと。 \\
API & アプリケーションプログラミングインタフェースのことである。ライブラリやサービスを利用するために公開された操作、引数、意味の取り決めである。 \\
\bottomrule
\end{longtable}
